\scanpage{9}
\mspage{4}
% archive foliation: 8
fixes qui porte la structure\ldots).
On peut aussi décrire ces structures en termes de $S$ ou de $A$ par la donnée
d'un groupe $\unsure{R|\Gamma}$ (``forme'' de $\mathbb{Z}/4\mathbb{Z}$) opérant
sur $S$ resp. sur $A$, qui devient un torseur sous $\Gamma$
\[
\begin{array}{ll}
(2\ \mathrm{bis}) & S\text{ torseur sous }\Gamma\\
(3\ \mathrm{bis}) & A\text{ torseur sous }\Gamma
\end{array}
\]
ce qui fait finalement \deleted{quatre}\inserted{cinq} descriptions de
\deleted{\unsure{la}} $=$ structure de ``carré''.
Je vais expliciter les passages de l'une de ces descriptions à l'autre.

\diagram{\begin{array}{rcl}
2\ \mathrm{bis})&\mapsto&2)\\
3\ \mathrm{bis})&\mapsto&3)
\end{array}}
en prenant l'unique élément $\underline{a}$ d'ordre 2 de $\Gamma$,
qui est une involution ss pts fixes de $S$ (resp $A$), et on pose
\[
D=S/\underline{a},\qquad\text{resp. }C=A/\underline{a}.
\]

$2)\deleted{\mathrm{bis}}\Longrightarrow1)$
Soit $A\subset\mathcal{P}_2(S)\inserted{\subset\mathcal{P}(S)}$ défini par
$A=\mathcal{P}_2(S)\setminus\operatorname{Im}D$,
c'est un ens à 4 éléments, \deleted{qui} et on a une relation tautologique
$R\subset S\times A$.

$3)\Longrightarrow1)$
Soit $S\subset\mathcal{P}_2(A)$ défini par
$S=\mathcal{P}_2(A)\setminus\operatorname{Im}C$,
ens à 4 éléments et on a relation tautologique $R'\subset A\times S$ d'où
$\unsure{{}^{s}R'|{}^{t}R'}=\unsure{{}^{\mathrm{def}}}R\subset S\times A$.

\diagram{\begin{array}{rcl}
1&\Longrightarrow&2\ \mathrm{bis}\\
1&\longrightarrow&3\ \mathrm{bis}
\end{array}}
On considère $\mathfrak{D}=\operatorname{Aut}\underline{Q}$
\margin{ordre 8\qquad $Q=\{S,A,R\}$}
\deleted{\unsure{et}}\deleted{$\unsure{\mathfrak{D}\cap V_Q}$}
\inserted{2-sous-gr. de Sylow de $\unsure{\mathcal{G}_S|\mathfrak{S}_S}$ ou
$\unsure{\mathcal{G}_A|\mathfrak{S}_A}$}
l'image inverse de $V_S\subset\unsure{\mathcal{G}_S|\mathfrak{S}_S}$ et celle de
$V_A\subset\unsure{\mathcal{G}_A|\mathfrak{S}_A}$ sont deux sous-groupes d'ordre 4
\deleted{\unsure{distincts}} de $\mathfrak{D}$, notés
$\unsure{V_S}$ et $V_A$, isomorphes respectivement à
$\mathbb{F}_2^C$ et $\mathbb{F}_2^D$.

\margin{\diagram{\begin{tikzcd}
& \unsure{\mathcal{G}_S|\mathfrak{S}_S}\\
\deleted{\unsure{G_Q}}\arrow[ur]\arrow[dr]&\\
&\unsure{\mathcal{G}_A|\mathfrak{S}_A}
\end{tikzcd}}}
\margin{\diagram{\begin{tikzcd}
\mathbb{F}_2^C\simeq V_S\arrow[r,hook]\arrow[d]&
\unsure{\mathcal{G}_S|\mathfrak{S}_S}\\
\mathfrak{D}\arrow[ur]\arrow[dr]&\\
\mathbb{F}_2^D\simeq V_A\arrow[u]\arrow[r,hook]&
\unsure{\mathcal{G}_A|\mathfrak{S}_A}
\end{tikzcd}}}
\deleted{(\illegible{1} $\unsure{A}$ \illegible{3} $\unsure{\mathfrak D}$ \illegible{2})
--- leur intersection est leur sous-groupe
$\unsure{V_Q}$ (\unsure{de rg $2$ sur $\mathbb{F}_2$}),
\unsure{à deux|il ne reste que deux} éléments d'ordre $4$
inverses l'un de l'autre, engendrant}
